\begin{center}
\begin{minipage}{\linewidth}
\centering
\resizebox{0.96\linewidth}{!}{%
\begin{tikzpicture}[
  >=Latex,
  every node/.style={font=\small},
  flow/.style={-{Latex[length=2.3mm]},thick,draw=GrisOscuro},
  point/.style={circle,fill=GrisOscuro,inner sep=1.8pt}
]
  \begin{scope}
    \clip[rounded corners=7pt] (0.2,0.25) rectangle (13.8,6.0);
    \fill[VerdeClaro] (0.2,0.25) rectangle (13.8,6.0);
    \fill[AzulClaro]
      (0.2,0.25) -- (6.35,0.25)
      .. controls (5.55,1.15) and (6.85,1.75) .. (6.05,2.55)
      .. controls (5.30,3.25) and (6.65,4.25) .. (6.35,6.0)
      -- (0.2,6.0) -- cycle;
  \end{scope}
  \draw[rounded corners=7pt,very thick,GrisOscuro] (0.2,0.25) rectangle (13.8,6.0);
  \node[anchor=north west,font=\bfseries] at (0.45,5.82) {Espacio de fases \(M\)};
  \node[font=\bfseries,Azul] at (2.55,5.15) {$\mathcal B(A_1)$};
  \node[font=\bfseries,Verde] at (10.55,5.15) {$\mathcal B(A_h)$};

  % Frontera común, punto silla y transporte sobre la frontera.
  \draw[very thick,Naranja]
    (6.35,0.25)
    .. controls (5.55,1.15) and (6.85,1.75) .. (6.05,2.55)
    .. controls (5.30,3.25) and (6.65,4.25) .. (6.35,6.0);
  \draw[very thick,GrisOscuro,densely dashed]
    (6.00,2.30) .. controls (5.55,3.05) and (6.55,3.75) .. (6.24,4.45);
  \node[point,label=right:$S$] (s) at (6.00,3.05) {};
  \node[align=left,fill=white,inner sep=2pt] at (7.78,4.15)
    {componente candidata\\$W^s(S)$};
  \node[point,label=left:$x$] (xb0) at (6.20,1.15) {};
  \node[point,label=right:$\varphi_t(x)$] (xb1) at (6.08,2.35) {};
  \draw[flow] (xb0) .. controls (5.78,1.62) and (6.35,1.95) .. (xb1);
  \node[rotate=90,Naranja,font=\bfseries] at (5.45,4.90)
    {$\partial\mathcal B(A_h)$};

  % Atractores y vecindad de equilibrio que no toca la cuenca oculta.
  \draw[very thick,Azul] (2.55,2.55) circle (0.43);
  \fill[Azul] (2.55,2.55) circle (2pt);
  \node[below=2pt] at (2.55,2.05) {$A_1$};
  \draw[very thick,Verde] (10.65,2.70) ellipse (0.78 and 0.42);
  \node[below=5pt] at (10.65,2.15) {$A_h$ oculto};
  \draw[flow,draw=Azul] (4.40,1.10) .. controls (3.55,1.45) and (3.0,1.90) .. (2.75,2.30);
  \draw[flow,draw=Verde] (12.75,4.25) .. controls (11.95,3.75) and (11.35,3.30) .. (10.95,2.98);

  \fill[GrisOscuro] (1.50,4.00) circle (2pt);
  \node[left=2pt] at (1.50,4.00) {$e$};
  \draw[GrisOscuro,dashed,thick] (1.50,4.00) circle (0.72);
  \node[above] at (1.50,4.72) {$U_e$};
  \node[draw=GrisOscuro,rounded corners,fill=white,inner sep=3pt]
    at (3.28,3.95) {$U_e\cap\mathcal B(A_h)=\varnothing$};
\end{tikzpicture}%
}

\smallskip
\begin{minipage}{0.92\linewidth}
\small\textbf{Lectura pedagógica.}
La ocultedad se expresa por vecindades de los equilibrios que no intersectan la
cuenca del atractor oculto. La frontera de una cuenca es invariante bajo un
flujo completo, pero identificarla globalmente con $W^s(S)$ exige excluir
otros conjuntos invariantes; el trazo discontinuo representa sólo una
componente candidata.
\end{minipage}
\end{minipage}
\end{center}
